\documentclass{article} % For LaTeX2e
\usepackage{iclr2021_conference,times}
\usepackage[hyphens]{url}
\usepackage{graphicx}

% Optional math commands from https://github.com/goodfeli/dlbook_notation.
\input{math_commands.tex}

\usepackage{hyperref}
\usepackage{algorithm}
\usepackage{algorithmic}
% \newcommand{\theHalgorithm}{\arabic{algorithm}}
% \newcommand{\theHalgorithm}{\arabic{algorithm}}

\title{SPG-Bandit: Skill Profile Guided Bandit for Task Selection in LLM Agent Skill Evolution}

%\iclrfinalcopy

\author{Author Name \\
Affiliation \\
Address \\
\texttt{email} \\
}

\begin{document}

\maketitle

\section{Introduction}
\label{sec:introduction}

\subsection{Background and Motivation}

Large language model (LLM) agents have demonstrated competence on complex tasks in diverse environments, including embodied control, web navigation, and software engineering. A key mechanism supporting this breadth of capability is \emph{skill evolution}, through which an agent reflects on execution trajectories, extracts reusable knowledge, and accumulates it in a skill library for later use.

Existing work has primarily focused on two components of this process. The first is \emph{skill quality}, namely the extraction of informative and reusable knowledge from trajectories. For example, SkillRL uses reflection to identify reusable strategies from successful episodes and recurrent failure modes from unsuccessful episodes. The second is \emph{skill utilization}, namely the organization, retrieval, and incorporation of skills into an agent's inference-time prompt. Representative approaches include structured skill banks with category-specific retrieval, vector database retrieval, and dynamic prompting.

A critical dimension remains comparatively underexplored: \textbf{the order in which tasks are selected for practice.} Existing systems commonly cycle uniformly through a fixed task pool, sample tasks at random, or rely on hand-designed curricula that do not adapt to an agent's evolving competence. Since tasks can make substantially different contributions to skill progress, an adaptive selection policy may use the interaction budget more efficiently. This motivates the following question: \emph{Can task selection conditioned on an agent's current skill profile improve the efficiency and effectiveness of skill evolution in LLM agents?}

\subsection{Contributions}

We propose \textbf{SPG-Bandit} (Skill Profile Guided Bandit), a framework that casts task selection as a contextual bandit problem informed by an agent's multidimensional skill profile.

The central idea is to use Multidimensional Item Response Theory (MIRT) to estimate latent proficiency across \(K\) skill dimensions and then select tasks expected to improve dimensions with lower estimated proficiency.

Specifically, our contributions are threefold:
\begin{enumerate}
    \item \textbf{Problem formulation.} We formalize task selection during LLM agent skill evolution as a nonstationary contextual bandit problem. The formulation captures both task specific differences in expected skill progress and the multidimensional nature of agent competence.
    \item \textbf{Method.} We propose SPG-Bandit, which combines MIRT based latent skill profiling with a gap weighted sliding window UCB policy. The method maintains an online proficiency estimate and selects tasks with high predicted progress on underdeveloped dimensions while accounting for estimation uncertainty.
    \item \textbf{General framework.} We provide a modular framework that can be integrated with existing skill-evolution methods and supports systematic evaluation of task-selection efficiency and downstream task performance.
\end{enumerate}

\section{Problem Formulation}
\label{sec:problem}

\subsection{Notation and Setup}

Consider an LLM agent interacting with a fixed task pool \(\mathcal{T} = \{\tau_1, \tau_2, \dots, \tau_M\}\) containing \(M\) tasks. At round \(t=1,2,\dots\), the agent selects and executes one task \(\tau_t\in\mathcal{T}\) and observes a binary outcome \(y_t\in\{0,1\}\), where one denotes success and zero denotes failure.

The agent has an unobserved \emph{true latent skill profile} \(\mathbf{s}_t^\star\in[0,1]^K\), which represents its proficiency across \(K\) skill dimensions. The algorithm instead uses the MIRT posterior estimate \(\widehat{\mathbf{s}}_t\in[0,1]^K\), which is measurable with respect to the history available before round \(t\). The true skill change is \(\boldsymbol{\Delta}_t^\star=\mathbf{s}_{t+1}^\star-\mathbf{s}_t^\star\). Because it is unobserved, the bandit receives the surrogate progress observation \(\widehat{\boldsymbol{\delta}}_t=\widehat{\mathbf{s}}_{t+1}-\widehat{\mathbf{s}}_t\). Thus, the method optimizes estimated progress, which is the quantity available from the MIRT filter, rather than directly optimizing \(\boldsymbol{\Delta}_t^\star\).

Each task \(\tau\) has a \emph{task embedding} \(\mathbf{e}_\tau\in\mathbb{R}^{d_e}\), where \(d_e\in\mathbb{N}_{+}\) is the embedding dimension. The embedding can be obtained by applying a sentence encoder to the natural language description of the task.

\subsection{Problem Definition}

We formalize adaptive task selection as a \emph{nonstationary contextual bandit} problem. At each round \(t\):

\begin{enumerate}
    \item The agent observes its context: the estimated profile \(\widehat{\mathbf{s}}_t\) and the embedding \(\mathbf{e}_\tau\) of every task \(\tau \in \mathcal{T}\).
    \item The agent selects an arm \(\tau_t \in \mathcal{T}\).
    \item The agent executes \(\tau_t\), observes the binary outcome \(y_t\), and updates its skills through reflection.
    \item The true profile evolves to \(\mathbf{s}_{t+1}^\star\), and MIRT produces the next estimate \(\widehat{\mathbf{s}}_{t+1}\).
\end{enumerate}

% To prioritize underdeveloped skill dimensions, 
To prioritize dimensions with lower estimated proficiency, define the \emph{gap weight} vector
\[
    \widehat{\mathbf{g}}_t = \operatorname{softmax}\!\left( \frac{\mathbf{1} - \widehat{\mathbf{s}}_t}{\gamma_g} \right),
\]

where \(\mathbf{1}\in\mathbb{R}^K\) is the all ones vector and \(\gamma_g>0\) is a temperature parameter. We use gap weighted skill progress as the immediate bandit reward,

\[
    \widehat r_t = \widehat{\mathbf{g}}_t^\top\widehat{\boldsymbol{\delta}}_t.
\]

The objective is to maximize the expected cumulative reward over a bandit horizon of \(T\) rounds,

\[
    \max_{\pi}\ \mathbb{E}\left[\sum_{t=1}^{T} \widehat r_t\ \middle|\ \tau_t\sim\pi(\cdot\mid\widehat{\mathbf{s}}_t,\{\mathbf{e}_\tau\})\right].
\]

where \(\pi\) denotes the task-selection policy and \(T\) is the number of rounds in the bandit phase. This reward prioritizes progress on underdeveloped dimensions without requiring an auxiliary objective.

We learn a predictor \(\widehat{\boldsymbol{\delta}}_t(\tau)\) of the surrogate progress and use \(\widehat{\mathbf{g}}_t^\top\widehat{\boldsymbol{\delta}}_t(\tau)\) as the estimated bandit reward. A superscript \(\star\) denotes an unobserved true quantity or an oracle regression parameter, whereas a hat denotes an estimate computed from data. Fixed task embeddings and algorithmic tuning parameters carry neither marker.


\section{Method}
\label{sec:method}

\subsection{Overview}

SPG-Bandit consists of a warmup phase and a bandit phase. During \textbf{warmup}, the agent samples tasks uniformly to obtain initial interaction outcomes. A sequential MIRT model estimates the latent profile sequence \(\{\widehat{\mathbf{s}}_t\}\), the surrogate progress vectors \(\{\widehat{\boldsymbol{\delta}}_t\}\), and task item parameters. These estimates supervise an MLP that maps a task embedding and the current estimated profile to a representation suitable for predicting surrogate progress.
    
During the \textbf{bandit phase}, the method computes \(\widehat{\mathbf{g}}_t\) from \(\widehat{\mathbf{s}}_t\) and evaluates every candidate task using a sliding window neural linear UCB rule. The MLP supplies a fixed representation of the observed estimated context, while the ridge head estimates gap weighted surrogate progress and its uncertainty. After execution, an online MIRT update produces \(\widehat{\mathbf{s}}_{t+1}\), from which \(\widehat{\boldsymbol{\delta}}_t\) is computed. The tuple \((\boldsymbol{\phi}_{t,\tau_t},\widehat{\boldsymbol{\delta}}_t)\) is then added to the sliding window.

\subsection{Key Components}

\paragraph{Multidimensional IRT (MIRT).}
The MIRT data generating model associates each task \(\tau\) with an unobserved discrimination vector \(\mathbf{a}_\tau^\star\in\mathbb{R}^K\) and an unobserved difficulty \(d_\tau^\star\in\mathbb{R}\). Given \(\mathbf{s}_t^\star\), its success probability is

\[
    p_\tau^\star(\mathbf{s}_t^\star)=\sigma\!\left(\mathbf{a}_\tau^{\star\top}\mathbf{s}_t^\star-d_\tau^\star\right).
\]

where \(\sigma(x)=1/(1+e^{-x})\). The estimator constrains \(\widehat{\mathbf{a}}_\tau\geq0\) elementwise and \(\widehat d_\tau\geq0\). The first constraint ensures that increasing an estimated skill component cannot reduce the fitted success probability. These constraints provide a directionally interpretable skill scale.

After collecting \(N\) warmup outcomes \(\{(\tau_t,y_t)\}_{t=1}^N\), we estimate profiles and item parameters by alternating optimization. The estimator uses \(p(\widehat{\mathbf{s}}_{t+1}\mid\widehat{\mathbf{s}}_t)=\mathcal{N}_{[0,1]^K}(\widehat{\mathbf{s}}_t,\boldsymbol{\Sigma}_s)\), where \(\boldsymbol{\Sigma}_s\succ0\) is a fixed temporal covariance matrix and \(\mathcal{N}_{[0,1]^K}\) denotes a Gaussian distribution truncated to \([0,1]^K\). This prior regularizes adjacent estimates and stabilizes inference.

\textit{E-step.} With \(\{(\widehat{\mathbf{a}}_\tau,\widehat d_\tau)\}\) fixed, compute \(\widehat{\mathbf{s}}_{1:N+1}\) by maximizing the joint log posterior
\[
\log p(\widehat{\mathbf{s}}_1) + \sum_{t=1}^{N}\left[\log p(y_t\mid\widehat{\mathbf{s}}_t,\widehat{\mathbf{a}}_{\tau_t},\widehat d_{\tau_t}) + \log p(\widehat{\mathbf{s}}_{t+1}\mid\widehat{\mathbf{s}}_t)\right],
\]
where \(p(\mathbf{s}_1)\) is a truncated Gaussian prior on \([0,1]^K\) centered at \(0.5\).

\textit{M-step.} With \(\{\widehat{\mathbf{s}}_t\}\) fixed, update \((\widehat{\mathbf{a}}_\tau,\widehat d_\tau)\) using L-BFGS-B, subject to \(\widehat{\mathbf{a}}_\tau\geq0\), \(\widehat d_\tau\geq0\), and \(\ell_2\) regularization on \(\widehat{\mathbf{a}}_\tau\).

We alternate the E-step and M-step until the objective converges.

\textit{Online update.} During the bandit phase, each observation \((y_t,\tau_t)\) refines the profile estimate without rerunning the full alternating procedure:

\[
    \widehat{\mathbf{s}}_{t+1}=\arg\max_{\mathbf{s}\in[0,1]^K}\left\{\log p(y_t\mid\mathbf{s},\widehat{\mathbf{a}}_{\tau_t},\widehat d_{\tau_t})+\log p(\mathbf{s}\mid\widehat{\mathbf{s}}_t)\right\}.
\]



\paragraph{Delta prediction via MLP.}
To predict surrogate progress, we train a two layer MLP \(\widehat\Phi\). Let \(\operatorname{cat}(\mathbf{u},\mathbf{v})\) denote vector concatenation. The learned encoder maps \(\operatorname{cat}(\mathbf{e}_\tau,\widehat{\mathbf{s}})\) to \(\boldsymbol{\phi}\in\mathbb{R}^{d_f}\). Its parameters are estimates and are denoted by \(\widehat{\mathbf{W}}_1,\widehat{\mathbf{b}}_1,\widehat{\mathbf{W}}_2,\widehat{\mathbf{b}}_2\).

\begin{align*}
    \mathbf{h} &= \max\!\left(\widehat{\mathbf{W}}_1\operatorname{cat}(\mathbf{e}_\tau,\widehat{\mathbf{s}})+\widehat{\mathbf{b}}_1,0\right), \\
    \boldsymbol{\phi} &= \widehat{\mathbf{W}}_2 \mathbf{h} + \widehat{\mathbf{b}}_2.
\end{align*}

A linear pretraining head \(\widehat{\mathbf{W}}_3\in\mathbb{R}^{d_f\times K}\) predicts surrogate progress as \(\widehat{\boldsymbol{\delta}}=\widehat{\mathbf{W}}_3^\top\boldsymbol{\phi}\). The encoder is trained on \(\{(\operatorname{cat}(\mathbf{e}_{\tau_i},\widehat{\mathbf{s}}_i),\widehat{\boldsymbol{\delta}}_i)\}_{i=1}^N\). After pretraining, \(\widehat{\mathbf{W}}_3\) is discarded and \(\widehat\Phi\) is fixed.

\paragraph{Gap-weighted UCB selection.}
To accommodate a non stationary context to progress relationship, we maintain a sliding window containing the \(\ell\) most recent \emph{completed} interactions and fit a local ridge head within that window. Let \(\mathcal{W}_t\) denote the indices stored in the buffer immediately before bandit round \(t\). The buffer is initialized with the final \(\ell\) warmup observations. For each buffered observation \(i\), define \(\boldsymbol{\phi}_i=\widehat\Phi(\operatorname{cat}(\mathbf{e}_{\tau_i},\widehat{\mathbf{s}}_i))\). With ridge regularization parameter \(\lambda>0\), the covariance and response matrices are

\[
    \widehat{\mathbf{W}}_t = \mathbf{A}_t^{-1}\mathbf{B}_t,
    \quad
    \mathbf{A}_t = \lambda \mathbf{I}_{d_f} + \sum_{i \in \mathcal{W}_t} \boldsymbol{\phi}_{i}\boldsymbol{\phi}_{i}^\top,
    \quad
    \mathbf{B}_t = \sum_{i \in \mathcal{W}_t} \boldsymbol{\phi}_{i} \widehat{\boldsymbol{\delta}}_{i}^\top,
\]

For a candidate task \(\tau\), let \(\boldsymbol{\phi}_{t,\tau}=\widehat\Phi(\operatorname{cat}(\mathbf{e}_\tau,\widehat{\mathbf{s}}_t))\). The predicted surrogate progress is \(\widehat{\boldsymbol{\delta}}_t(\tau)=\widehat{\mathbf{W}}_t^\top\boldsymbol{\phi}_{t,\tau}\). The UCB score is

\[
    \operatorname{score}_t(\tau) = \widehat{\mathbf{g}}_t^\top \widehat{\boldsymbol{\delta}}_t(\tau) + \beta_t\lVert\widehat{\mathbf{g}}_t\rVert_2\sqrt{\boldsymbol{\phi}_{t,\tau}^\top\mathbf{A}_t^{-1}\boldsymbol{\phi}_{t,\tau}},
\]

The first term favors tasks with high predicted gap weighted progress. The second term is a confidence bonus for the scalarized reward \(\mathbf{g}_t^\top\boldsymbol{\delta}_t\), where \(\beta_t>0\) is selected according to a ridge regression confidence bound.

\paragraph{Sliding-window update.}
After observing \(\boldsymbol{\delta}_t\), we append \((\boldsymbol{\phi}_{t,\tau_t},\boldsymbol{\delta}_t)\) to the window and evict its oldest entry when the window exceeds size \(\ell\). The incremental updates are:

\begin{align*}
    \mathbf{A}_{t+1} &= \mathbf{A}_t + \boldsymbol{\phi}_{t,\tau_t}\boldsymbol{\phi}_{t,\tau_t}^\top - \boldsymbol{\phi}_{t-\ell}\boldsymbol{\phi}_{t-\ell}^\top, \\
    \mathbf{B}_{t+1} &= \mathbf{B}_t + \boldsymbol{\phi}_{t,\tau_t}\boldsymbol{\delta}_t^\top - \boldsymbol{\phi}_{t-\ell}\boldsymbol{\delta}_{t-\ell}^\top, \\
    \widehat{\mathbf{W}}_{t+1} &= \mathbf{A}_{t+1}^{-1}\mathbf{B}_{t+1}.
\end{align*}

The window size \(\ell\) controls the tradeoff between estimation quality, which benefits from a larger sample, and responsiveness to nonstationarity, which benefits from recent observations.

\subsection{Algorithm}

Algorithm~\ref{alg:spg} summarizes the complete procedure.

\begin{algorithm}[t]
\caption{SPG-Bandit}
\label{alg:spg}
\begin{algorithmic}
\STATE {\bf Input:} Task pool \(\mathcal{T}\), warmup steps \(N\geq\ell\), bandit horizon \(T\), window size \(\ell\), feature dimension \(d_f\), gap temperature \(\gamma_g\), confidence radii \(\{\beta_t\}\), regularizer \(\lambda>0\), skill dimensions \(K\)
\STATE {\bf Phase 1: Warmup}
\FOR{\(i = 1\) to \(N\)}
    \STATE Select \(\tau_i \in \mathcal{T}\) uniformly
    \STATE Execute \(\tau_i\), observe \(y_i\in\{0,1\}\), and perform the skill-evolution update
    \STATE Store \((\tau_i,y_i)\)
\ENDFOR
\STATE Run sequential MIRT EM to obtain \(\{\widehat{\mathbf{s}}_i\}_{i=1}^{N+1}\), \(\{\widehat{\boldsymbol{\delta}}_i\}_{i=1}^{N}\), and \(\{\widehat{\mathbf{a}}_\tau,\widehat d_\tau\}_{\tau\in\mathcal{T}}\)
\STATE Train and freeze \(\widehat\Phi\) on \(\{(\operatorname{cat}(\mathbf{e}_{\tau_i},\widehat{\mathbf{s}}_i),\widehat{\boldsymbol{\delta}}_i)\}_{i=1}^{N}\)
\STATE Reindex the final \(\ell\) warmup tuples as \(\{(\tau_i,\mathbf{s}_i,\boldsymbol{\phi}_i,\boldsymbol{\delta}_i)\}_{i=1-\ell}^{0}\), where \(\boldsymbol{\phi}_i=\Phi(\operatorname{cat}(\mathbf{e}_{\tau_i},\mathbf{s}_i))\)
\STATE Set \(\mathcal{W}_1=\{1-\ell,\ldots,0\}\), \(\mathbf{A}_1=\lambda\mathbf{I}_{d_f}+\sum_{i\in\mathcal{W}_1}\boldsymbol{\phi}_i\boldsymbol{\phi}_i^\top\), \(\mathbf{B}_1=\sum_{i\in\mathcal{W}_1}\boldsymbol{\phi}_i\widehat{\boldsymbol{\delta}}_i^\top\), \(\widehat{\mathbf{W}}_1=\mathbf{A}_1^{-1}\mathbf{B}_1\), and \(\widehat{\mathbf{s}}_1=\widehat{\mathbf{s}}_{N+1}\)
\STATE {\bf Phase 2: Bandit}
\FOR{\(t = 1\) to \(T\)}
    \STATE Compute \(\widehat{\mathbf{g}}_t=\operatorname{softmax}((\mathbf{1}-\widehat{\mathbf{s}}_t)/\gamma_g)\)
    \FOR{each \(\tau \in \mathcal{T}\)}
        \STATE \(\boldsymbol{\phi}_{t,\tau}=\widehat\Phi(\operatorname{cat}(\mathbf{e}_\tau,\widehat{\mathbf{s}}_t))\), \(\widehat{\boldsymbol{\delta}}_t(\tau)=\widehat{\mathbf{W}}_t^\top\boldsymbol{\phi}_{t,\tau}\)
        \STATE \(\operatorname{score}_t(\tau)=\widehat{\mathbf{g}}_t^\top\widehat{\boldsymbol{\delta}}_t(\tau)+\beta_t\lVert\widehat{\mathbf{g}}_t\rVert_2\sqrt{\boldsymbol{\phi}_{t,\tau}^\top\mathbf{A}_t^{-1}\boldsymbol{\phi}_{t,\tau}}\)
    \ENDFOR
    \STATE \(\tau_t=\arg\max_{\tau\in\mathcal{T}}\operatorname{score}_t(\tau)\)
    \STATE Execute \(\tau_t\) and observe \(y_t\)
    \STATE \(\widehat{\mathbf{s}}_{t+1}\gets\textsc{BayesUpdate}(\widehat{\mathbf{s}}_t,\widehat{\mathbf{a}}_{\tau_t},\widehat d_{\tau_t},y_t)\)
    \STATE \(\widehat{\boldsymbol{\delta}}_t\gets\widehat{\mathbf{s}}_{t+1}-\widehat{\mathbf{s}}_t\) and \(\boldsymbol{\phi}_t\gets\boldsymbol{\phi}_{t,\tau_t}\)
    \STATE \(\mathcal{W}_{t+1}\gets(\mathcal{W}_t\cup\{t\})\setminus\{t-\ell\}\)
    \STATE \(\mathbf{A}_{t+1}\gets\mathbf{A}_t+\boldsymbol{\phi}_t\boldsymbol{\phi}_t^\top-\boldsymbol{\phi}_{t-\ell}\boldsymbol{\phi}_{t-\ell}^\top\)
    \STATE \(\mathbf{B}_{t+1}\gets\mathbf{B}_t+\boldsymbol{\phi}_t\widehat{\boldsymbol{\delta}}_t^\top-\boldsymbol{\phi}_{t-\ell}\widehat{\boldsymbol{\delta}}_{t-\ell}^\top\)
    \STATE \(\widehat{\mathbf{W}}_{t+1}\gets\mathbf{A}_{t+1}^{-1}\mathbf{B}_{t+1}\)
\ENDFOR
\end{algorithmic}
\end{algorithm}

\section{Theoretical Analysis}
\label{sec:theory}

Because the relationship between context and skill progress may change over time, expected task rewards can be non stationary. In this section, \(t\in\{1,\ldots,T\}\) indexes only bandit phase rounds. We analyze the dynamic \emph{pseudo regret} of SPG-Bandit under the sliding window ridge estimator described in Section~\ref{sec:method}.

\subsection{Non-stationary Setting}

For each candidate task \(\tau\), define \(\boldsymbol{\phi}_{t,\tau}=\widehat\Phi(\operatorname{cat}(\mathbf{e}_\tau,\widehat{\mathbf{s}}_t))\) and the conditional mean surrogate progress
\[
\boldsymbol{\mu}_t(\tau)=\mathbb{E}[\widehat{\boldsymbol{\delta}}_t(\tau)\mid\mathcal{F}_{t-1}].
\]
The expected scalar reward is \(q_t(\tau)=\widehat{\mathbf{g}}_t^\top\boldsymbol{\mu}_t(\tau)\). Let \(\tau_t^*\in\arg\max_{\tau\in\mathcal{T}}q_t(\tau)\) be the dynamic oracle action. We evaluate the dynamic pseudoregret

\[
    R_T=\sum_{t=1}^{T}\bigl[q_t(\tau_t^*)-q_t(\tau_t)\bigr].
\]

This quantity is a pseudo regret because it compares conditional expected rewards rather than realized noisy outcomes. Let \(\mathcal{F}_{t-1}\) denote the history available before round \(t\). The estimated context \(\widehat{\mathbf{s}}_t\), and hence \(\widehat{\mathbf{g}}_t\) and \(\boldsymbol{\phi}_{t,\tau}\), are \(\mathcal{F}_{t-1}\)-measurable.

\subsection{Assumptions}

\begin{enumerate}
    \item \textbf{(Time varying linear surrogate model.)} There are oracle matrices \(\{\mathbf{W}_t^*\}_{t=1}^{T}\), where \(\mathbf{W}_t^*\in\mathbb{R}^{d_f\times K}\), such that \(\widehat{\boldsymbol{\delta}}_t(\tau)=\mathbf{W}_t^{*\top}\boldsymbol{\phi}_{t,\tau}+\boldsymbol{\varepsilon}_t(\tau)\), with \(\mathbb{E}[\boldsymbol{\varepsilon}_t(\tau)\mid\mathcal{F}_{t-1},\tau]=\mathbf{0}\).
    \item \textbf{(Boundedness and noise.)} \(\lVert\boldsymbol{\phi}_{t,\tau}\rVert_2\leq1\), \(\lVert\mathbf{W}_t^*\rVert_F\leq S\), and, for every \(\mathbf{u}\) satisfying \(\lVert\mathbf{u}\rVert_2\leq1\), the scalar \(\mathbf{u}^\top\boldsymbol{\varepsilon}_t(\tau)\) is conditionally \(\sigma\)-sub-Gaussian.
    \item \textbf{(Parameter path variation.)} The reward model variation is \(P_T=\sum_{t=1}^{T-1}\lVert\mathbf{W}_{t+1}^*-\mathbf{W}_t^*\rVert_F\).
    \item \textbf{(Calibrated window and coverage.)} Before the first bandit round, warmup provides \(\ell\) calibration observations generated under \(\mathbf{W}_1^*\). At each bandit round, the window contains exactly \(\ell\) observations and satisfies \(\sum_{i\in\mathcal{W}_t}\boldsymbol{\phi}_i\boldsymbol{\phi}_i^\top\succeq\kappa\ell\mathbf{I}_{d_f}\) for a constant \(\kappa>0\).
\end{enumerate}

Assumption~1 matches the algorithmic regression target, which is \(\boldsymbol{\delta}_t\) rather than \(\mathbf{s}_{t+1}\). Assumption~2 controls the scale of features, parameters, and noise. Assumption~3 quantifies non stationarity through the parameter path. Assumption~4 isolates the effect of non stationarity from an ill conditioned design matrix. For a failure probability \(\rho\in(0,1)\), Appendix~\ref{app:proof} provides a detailed high probability proof.

\subsection{Dynamic Regret Decomposition}

Let \(w_t(\tau)=\sqrt{\boldsymbol{\phi}_{t,\tau}^\top\mathbf{A}_t^{-1}\boldsymbol{\phi}_{t,\tau}}\) and \(\widehat{\mathbf{W}}_t=\mathbf{A}_t^{-1}\mathbf{B}_t\). Under Assumptions~1 to 4, Appendix~\ref{app:proof} proves that, with probability at least \(1-\rho\),

\[
\left|\mathbf{g}_t^\top(\widehat{\mathbf{W}}_t-\mathbf{W}_t^*)^\top\boldsymbol{\phi}_{t,\tau}\right|
\leq \beta_t\lVert\mathbf{g}_t\rVert_2w_t(\tau)+D_t,
\]

where \(\beta_t=\tilde{O}(\sigma\sqrt{K d_f}+\sqrt{\lambda}S)\) is a confidence radius and \(D_t\) is the explicitly bounded bias from stale observations. Since the policy selects the task with the largest UCB score, approximate optimism yields
\[
q_t(\tau_t^*)-q_t(\tau_t)
\leq 2\beta_t\lVert\mathbf{g}_t\rVert_2
 w_t(\tau_t)+2D_t.
\]

Because \(\mathbf{g}_t\) is a softmax vector, \(\lVert\mathbf{g}_t\rVert_2\leq1\). Assumption~4 gives \(w_t(\tau)\leq(\lambda+\kappa\ell)^{-1/2}\) for every candidate task. The detailed counting argument in Appendix~\ref{app:proof} gives \(\sum_{t=1}^{T}D_t\leq \ell^2P_T/\sqrt{\lambda(\lambda+\kappa\ell)}\). Therefore, with \(\beta=\max_{t\leq T}\beta_t\),

\[
    R_T \leq \underbrace{\frac{2\beta T}{\sqrt{\lambda+\kappa\ell}}}_{\text{estimation}}+\underbrace{\frac{2\ell^2P_T}{\sqrt{\lambda(\lambda+\kappa\ell)}}}_{\text{drift}}.
\]

\subsection{Optimal Window and Regret Bound}

For \(P_T>0\), when \(\kappa\ell\geq\lambda\), balancing the two terms yields the following feasible integer window size:

\[
    \ell^*=\min\!\left\{T,\max\!\left\{\left\lceil\frac{\lambda}{\kappa}\right\rceil,\left\lceil\left(\frac{\beta T\sqrt{\lambda}}{P_T}\right)^{1/2}\right\rceil\right\}\right\},
\]

and the corresponding cumulative regret bound, up to universal constants:

\[
    R_T=O\!\left(\frac{(\beta T)^{3/4}P_T^{1/4}}{\kappa^{1/2}\lambda^{1/8}}\right).
\]

If \(P_T=0\), the drift term vanishes and \(\ell=T\) minimizes the displayed bound. If \(P_T=o(T)\) and \(\beta\), \(\kappa^{-1}\), and \(\lambda^{-1}\) grow at most polylogarithmically, then the bound is sublinear. For example, if \(P_T=O(T^\nu)\) for \(\nu<1\), then

\[
    R_T=\tilde{O}\!\left(T^{(3+\nu)/4}\right)=o(T).
\]








% \bibliography{refs}
% \bibliographystyle{iclr2021_conference}

\appendix
\section{Detailed Dynamic Pseudo-Regret Analysis}
\label{app:proof}

This appendix proves the bound in Section~\ref{sec:theory}. The result is conditional on the time-varying linear delta model and the window-coverage condition stated there. These assumptions are necessary: a changing latent profile alone does not establish a non-stationary bandit regret bound.

\subsection{Filtration and confidence event}

Let \(\mathcal{F}_{t-1}\) be the sigma-field generated by all observations and selected tasks before bandit round \(t\). For the selected task, abbreviate \(\boldsymbol{\phi}_{t}=\boldsymbol{\phi}_{t,\tau_t}\). Assumption~1 can be written as
\begin{equation}
    \boldsymbol{\delta}_{t}=\mathbf{W}_t^{*\top}\boldsymbol{\phi}_{t}+\boldsymbol{\varepsilon}_{t},
    \qquad
    \mathbb{E}[\boldsymbol{\varepsilon}_{t}\mid\mathcal{F}_{t-1},\tau_t]=\mathbf{0}.
    \label{eq:appendix-linear-model}
\end{equation}
For \(\mathbf{u}\in\mathbb{R}^{K}\) with \(\lVert\mathbf{u}\rVert_2\leq1\), the scalar noise \(\mathbf{u}^{\top}\boldsymbol{\varepsilon}_t\) is conditionally \(\sigma\)-sub-Gaussian.

Recall the ridge matrices formed from completed observations in \(\mathcal{W}_t\):
\begin{equation}
    \mathbf{A}_{t}=\lambda\mathbf{I}_{d_f}+\sum_{i\in\mathcal{W}_t}\boldsymbol{\phi}_{i}\boldsymbol{\phi}_{i}^{\top},
    \qquad
    \mathbf{B}_{t}=\sum_{i\in\mathcal{W}_t}\boldsymbol{\phi}_{i}\boldsymbol{\delta}_{i}^{\top},
    \qquad
    \widehat{\mathbf{W}}_t=\mathbf{A}_t^{-1}\mathbf{B}_t.
    \label{eq:appendix-ridge}
\end{equation}
For a failure probability \(\rho\in(0,1)\), define
\begin{equation}
    \beta_t=\sigma\sqrt{2K\log\!\left(
    \frac{K\det(\mathbf{A}_t)^{1/2}}{\rho\lambda^{d_f/2}}
    \right)}+\sqrt{\lambda}S.
    \label{eq:appendix-beta}
\end{equation}
The factor \(K\) follows from applying a scalar self-normalized concentration inequality to each output coordinate and taking a union bound. It can be replaced by a sharper vector-valued confidence radius when additional noise structure is available.

\subsection{Estimation error and drift bias}

For a candidate \(\tau\), let
\[
    w_t(\tau)=\sqrt{\boldsymbol{\phi}_{t,\tau}^{\top}\mathbf{A}_t^{-1}\boldsymbol{\phi}_{t,\tau}}.
\]
The following lemma separates stochastic estimation error from bias due to stale samples.

\paragraph{Lemma 1 (windowed prediction error).}
Under Assumptions~1--4, with probability at least \(1-\rho\), simultaneously for every \(t\leq T\) and \(\tau\in\mathcal{T}\),
\begin{equation}
\left|
\mathbf{g}_t^{\top}(\widehat{\mathbf{W}}_t-\mathbf{W}_t^*)^{\top}
\boldsymbol{\phi}_{t,\tau}
\right|
\leq
\beta_t\lVert\mathbf{g}_t\rVert_2w_t(\tau)+D_t,
\label{eq:appendix-prediction-error}
\end{equation}
where, for a constant \(c_{\lambda,\kappa}>0\) depending only on the regularization and coverage constants,
\begin{equation}
    D_t\leq c_{\lambda,\kappa}
    \sum_{j=\max\{1,t-\ell\}}^{t-1}
    \lVert\mathbf{W}_{j+1}^{*}-\mathbf{W}_{j}^{*}\rVert_F.
    \label{eq:appendix-drift-bias}
\end{equation}

\paragraph{Proof.}
Substituting~\eqref{eq:appendix-linear-model} into~\eqref{eq:appendix-ridge} and adding and subtracting \(\mathbf{W}_t^*\) gives
\begin{align}
\widehat{\mathbf{W}}_t-\mathbf{W}_t^*
={}&\mathbf{A}_t^{-1}\sum_{i\in\mathcal{W}_t}
\boldsymbol{\phi}_i\boldsymbol{\varepsilon}_i^{\top}
-\lambda\mathbf{A}_t^{-1}\mathbf{W}_t^* \nonumber\\
&+\mathbf{A}_t^{-1}\sum_{i\in\mathcal{W}_t}
\boldsymbol{\phi}_i\boldsymbol{\phi}_i^{\top}
\left(\mathbf{W}_i^*-\mathbf{W}_t^*\right).
\label{eq:appendix-error-decomposition}
\end{align}
The first two terms are the usual stochastic and regularization terms of ridge regression. Applying the self-normalized inequality coordinatewise, followed by Cauchy--Schwarz over the \(K\) output coordinates, bounds their scalarized prediction error by
\[
\beta_t\lVert\mathbf{g}_t\rVert_2
\lVert\boldsymbol{\phi}_{t,\tau}\rVert_{\mathbf{A}_t^{-1}}.
\]
For the final term, Assumption~4 bounds the inverse window Gram matrix, while \(\lVert\boldsymbol{\phi}_i\rVert_2\leq1\). Thus, its scalarized prediction error is at most a constant \(c_{\lambda,\kappa}\) times
\[
\sum_{i\in\mathcal{W}_t}\lVert\mathbf{W}_i^*-\mathbf{W}_t^*\rVert_F.
\]
By telescoping \(\mathbf{W}_i^*-\mathbf{W}_t^*\) along the parameter path, this is bounded by the right-hand side of~\eqref{eq:appendix-drift-bias}. Combining the three bounds proves the lemma. \hfill\(\square\)

Summing~\eqref{eq:appendix-drift-bias} over time counts each increment \(\lVert\mathbf{W}_{j+1}^*-\mathbf{W}_j^*\rVert_F\) in at most \(\ell\) windows. Hence,
\begin{equation}
    \sum_{t=1}^{T}D_t\leq c_{\lambda,\kappa}\ell P_T.
    \label{eq:appendix-drift-sum}
\end{equation}

\subsection{From confidence to dynamic pseudo-regret}

Define the confidence width
\[
    U_t(\tau)=\beta_t\lVert\mathbf{g}_t\rVert_2w_t(\tau).
\]
On the event in Lemma~1, the expected reward satisfies, for every candidate \(\tau\),
\begin{equation}
    q_t(\tau)\leq
    \mathbf{g}_t^{\top}\widehat{\mathbf{W}}_t^{\top}\boldsymbol{\phi}_{t,\tau}
    +U_t(\tau)+D_t.
    \label{eq:appendix-upper-confidence}
\end{equation}
The selected task maximizes the first two terms on the right-hand side. Applying~\eqref{eq:appendix-upper-confidence} to \(\tau_t^*\), then applying the corresponding lower bound from Lemma~1 to \(\tau_t\), yields
\begin{align}
q_t(\tau_t^*)-q_t(\tau_t)
&\leq
\left[\mathbf{g}_t^{\top}\widehat{\mathbf{W}}_t^{\top}\boldsymbol{\phi}_{t,\tau_t}
+U_t(\tau_t)+D_t\right]-q_t(\tau_t) \nonumber\\
&\leq 2U_t(\tau_t)+2D_t.
\label{eq:appendix-instant-regret}
\end{align}
This is approximate optimism: the stale-sample term \(D_t\) prevents exact UCB optimism, but is explicitly controlled by the parameter-path variation.

Because \(\mathbf{g}_t\) is a softmax vector, \(\lVert\mathbf{g}_t\rVert_2\leq1\). Once the window is full, Assumption~4 gives \(\mathbf{A}_t\succeq\kappa\ell\mathbf{I}_{d_f}\). The warmup initialization supplies a full window before the first bandit decision. Consequently,
\begin{equation}
    \sum_{t=1}^{T}U_t(\tau_t)
    \leq\frac{\beta T}{\sqrt{\kappa\ell}},
    \qquad
    \beta=\max_{t\leq T}\beta_t.
    \label{eq:appendix-width-sum}
\end{equation}
Summing~\eqref{eq:appendix-instant-regret}, and using~\eqref{eq:appendix-drift-sum} and~\eqref{eq:appendix-width-sum}, proves
\[
R_T\leq
O\!\left(\frac{\beta T}{\sqrt{\kappa\ell}}\right)
+O(\ell P_T).
\]
Balancing the two terms gives the window choice and regret rate in Section~\ref{sec:theory}. \hfill\(\square\)


\end{document}
